\mychapter{Warianty algorytmu Dijkstry dla problemu miasta 15-minutowego}
\label{chap:warianty-dijkstry}

W~rozdziale~\ref{chap:sssp} uzasadniono wybór algorytmu Dijkstry jako podstawy
rozwiązania problemu miasta 15-minutowego. Klasyczna implementacja oparta na kopcu
binarnym, traktowana w~literaturze jako rozwiązanie domyślne, nie wykorzystuje
jednak w~pełni własności tego zadania: ograniczonego zasięgu izochron
wyznaczonych progiem $t$, wąskiego rozkładu wag krawędzi oraz powtarzalnej
struktury obliczeń po kategoriach usług.

Poniżej omówiono trzy warianty algorytmu Dijkstry, które posłużyły
do eksperymentów opisanych dalej w~pracy:
\begin{itemize}
  \item wariant \emph{bazowy}: klasyczny algorytm Dijkstry
        z~kopcem binarnym, bez ograniczenia progiem, w~którym warunek
        $\le t$ stosowany jest dopiero po zakończeniu przeszukiwania
        jako filtr na tablicy odległości,
  \item wariant \emph{ograniczony progiem} $t$: ten sam algorytm,
        w~którym przeszukiwanie kończy się w~chwili wyciągnięcia z~kopca
        wierzchołka o~odległości $>t$, dzięki czemu praca jest proporcjonalna
        do rozmiaru izochrony, a~nie do rozmiaru całego grafu,
  \item wariant z~\emph{kolejką kubełkową}: ten sam algorytm
        ograniczony progiem $t$, w~którym kopiec binarny zastępuje się
        jednopoziomową tablicą kubełków FIFO o~stałej szerokości $\Delta$,
        co~zmienia koszt amortyzowany na operację z~$O(\log|V|)$ na~$O(1)$.
\end{itemize}
Wspólnym elementem wszystkich wariantów jest \emph{redukcja wieloźródłowa}
oparta na konstrukcji superźródła~\cite{Lam2024}, omówiona w~podrozdziale
\ref{sec:warianty:supersource}.

\section{Redukcja wieloźródłowa: konstrukcja superźródła}
\label{sec:warianty:supersource}

Z~definicji zbioru dostępności
$C_k(t)=\{\,v\in V:\ \tau_k(v)\le t\,\}$,
gdzie $\tau_k(v)=\min_{s\in S_k} d_G(v,s)$
(podrozdz.~\ref{sec:intro:formalizacja}),
wynika, że dla kategorii usług $S_k$ poszukujemy minimum odległości
\emph{po wszystkich} źródłach $s\in S_k$. Naiwna realizacja wymagałaby
uruchomienia algorytmu SSSP osobno dla każdego $s\in S_k$, czyli $|S_k|$
przeszukiwań grafu. Dla rozważanych obszarów miejskich
liczność $|S_k|$ dochodzi w~niektórych kategoriach do kilkunastu tysięcy
(tab.~\ref{tab:uslugi_kategorie}), co czyni takie podejście nieefektywnym.

Redukcja wieloźródłowa zaproponowana przez Lama~\cite{Lam2024}
w~kontekście problemu miasta 15-minutowego sprowadza wariant
wieloźródłowy do jednorazowego uruchomienia SSSP. W~grafie $G$ dodaje się
\emph{wirtualny wierzchołek superźródłowy} $s^\star_k\notin V$ oraz krawędzie
o~wadze zerowej łączące go ze wszystkimi usługami kategorii:
\[
G^\star_k=\bigl(V\cup\{s^\star_k\},\ E\cup\{(s^\star_k,s):\,s\in S_k\}\bigr),
\qquad
w(s^\star_k,s)=0\text{ dla każdego }s\in S_k.
\]
Wówczas dla każdego $v\in V$ zachodzi równość
\[
d_{G^\star_k}(s^\star_k,v)=\min_{s\in S_k} d_G(s,v)=\tau_k(v),
\]
a~zbiór $C_k(t)$ wyznacza się jako
\[
C_k(t)=\{\,v\in V:\ d_{G^\star_k}(s^\star_k,v)\le t\,\}.
\]
Uzyskanie zbioru $C_k(t)$ wymaga zatem dokładnie jednego uruchomienia algorytmu
SSSP z~superźródła $s^\star_k$, niezależnie od liczności zbioru $S_k$.

Konstrukcja superźródła jest wspólna dla wszystkich trzech rozważanych
wariantów algorytmu, dzięki czemu różnice w~mierzonych czasach wykonania
można interpretować jako efekt wyłącznie zmiany algorytmu wewnętrznego,
a~nie różnicy w~sposobie obsługi wielu źródeł.

\clearpage

\section{Wariant bazowy: klasyczny algorytm Dijkstry}
\label{sec:warianty:regular}

Wariant bazowy jest klasycznym algorytmem Dijkstry z~kopcem
binarnym opisanym w~podrozdziale~\ref{sec:sssp:dijkstra}, uruchomionym
z~superźródła $s^\star_k$ na grafie $G^\star_k$. Przeszukiwanie nie jest
w~żaden sposób ograniczane progiem $t$: algorytm działa do wyczerpania
kolejki priorytetowej i~zwraca pełną tablicę odległości $d_{G^\star_k}(s^\star_k,v)$
dla wszystkich osiągalnych $v\in V$. Warunek dostępności w~czasie $\le t$
jest następnie nakładany jako filtr na uzyskaną tablicę odległości:
\[
C_k(t)=\{\,v\in V:\ d_{G^\star_k}(s^\star_k,v)\le t\,\}.
\]
Pełna ścieżka obliczeniowa wariantu bazowego polega na 
przeszukaniu całego osiągalnego grafu, a~następnie filtracji po progu.

\subsection{Rola wariantu bazowego w~porównaniu}
\label{subsec:warianty:regular:rola}

Wariant bazowy pełni tu podwójną rolę. Po pierwsze, służy jak wzorzec poprawności.
Przy ustalonym progu $t$ zwracane zbiory muszą
zgadzać się z~wynikami pozostałych dwóch wariantów. Weryfikacja polega na
porównaniu wyników dla tej samej pary (graf, kategoria) i~pozwala uchwycić
ewentualne błędy w~ograniczeniu progiem lub w~kolejce kubełkowej.

Po drugie, daje punkt odniesienia dla pomiarów czasu: różnica między
czasem wariantu bazowego a~czasem wariantu ograniczonego progiem~$t$
odzwierciedla wpływ samego ograniczenia zasięgu izochrony na koszt obliczeń,
czyli odpowiada na pytanie badawcze~\textbf{P1}.

\section{Wariant ograniczony progiem $t$}
\label{sec:warianty:modified}

Wariant ograniczony progiem~$t$ korzysta z~tej samej procedury co~wariant
bazowy, z~tym że przeszukiwanie kończy się w~momencie wyciągnięcia z~kopca
wierzchołka o~tymczasowej odległości większej niż~$t$. W~efekcie odwiedzana
jest w~zasadzie ta część grafu, która wpada w~izochronę, tj.\ zbiór
$\{v\in V:\,d_{G^\star_k}(s^\star_k,v)\le t\}\cup\{s^\star_k\}$.

\subsection{Schemat algorytmu}
\label{subsec:warianty:modified:schemat}

Schemat algorytmu różni się od wariantu bazowego w~dwóch miejscach: dodano
warunek wyjścia z~głównej pętli przy odległości większej od~$t$ oraz nie
wstawia się do kopca wierzchołków, których tymczasowa odległość przekraczałaby
ten próg.

\begin{codebox}
\Procname{$\proc{Bounded-Dijkstra}(G^\star_k, s^\star_k, t)$}
\li \For each $v\in V\cup\{s^\star_k\}$
\li \Do $\id{dist}[v]\gets\infty$ \End
\li $\id{dist}[s^\star_k]\gets 0$
\li $H\gets\{(0,s^\star_k)\}$
\li \While $H\neq\emptyset$
\li \Do $(d,u)\gets\proc{Extract-Min}(H)$
\li     \If $d > t$
\li     \Then \kw{break}                              \quad // przerwanie progiem $t$
        \End
\li     \If $d > \id{dist}[u]$
\li     \Then \kw{continue}
        \End
\li     \For each $(u,v)\in E(G^\star_k)$
\li     \Do $\id{nd}\gets d + w(u,v)$
\li         \If $\id{nd} > t$
\li         \Then \kw{continue}                       \quad // pruning relaksacji
            \End
\li         \If $\id{nd} < \id{dist}[v]$
\li         \Then $\id{dist}[v]\gets\id{nd}$
\li             $\proc{Heap-Insert}(H,(\id{nd},v))$
            \End
        \End
    \End
\li \Return $\{v:\,\id{dist}[v]\le t\}$
\end{codebox}

\subsection{Poprawność ograniczenia progiem}
\label{subsec:warianty:modified:poprawnosc}

Argument poprawności jest taki sam jak w~klasycznym algorytmie Dijkstry: gdy z~kopca
wyciągany jest wierzchołek $u$ o~odległości~$d$, żaden z~pozostających
w~kopcu elementów nie ma odległości mniejszej od~$d$.
Jeśli zatem $d>t$, to odległość każdego elementu nadal obecnego w~kopcu jest
co~najmniej $d$, a~więc większa niż~$t$; żaden wierzchołek osiągalny
w~czasie $\le t$ nie może już zostać dodany do izochrony. Przerwanie pętli
w~tym momencie jest zatem bezstratne.

Drugi warunek (pominięcie relaksacji prowadzącej do odległości $\id{nd}>t$)
jest analogicznym ograniczeniem przeszukiwania: wierzchołki o~tymczasowej odległości większej
od~$t$ nie wejdą do zwracanego zbioru $C_k(t)$, dlatego ich dalsze utrzymywanie
w~kopcu (a~tym bardziej rozszerzanie z~nich) jest zbędne. Pominięcie wstawienia
do kopca nie zmienia wyniku, ale zmniejsza rozmiar kopca w~trakcie
przeszukiwania, co~ma znaczenie praktyczne ze~względu na logarytmiczny
koszt operacji \proc{Heap-Insert} i~\proc{Extract-Min} w~funkcji liczności
kopca.

\clearpage

\subsection{Złożoność i oczekiwany zysk}
\label{subsec:warianty:modified:zlozonosc}

Asymptotycznie złożoność wariantu ograniczonego pozostaje
$\mathcal{O}((|V'|+|E'|)\log|V'|)$, gdzie $V'$ i~$E'$ oznaczają, odpowiednio,
zbiór wierzchołków i~krawędzi faktycznie odwiedzonych przez algorytm.
Zysk względem wariantu bazowego zależy zatem od~tego, jaki udział grafu
pokrywa izochrona $C_k(t)$: im mniejszy stosunek $|V'|/|V|$, tym więcej
pracy odcina próg.

W~problemie miasta 15-minutowego skala tego zysku jest jednak ograniczona
przez redukcję wieloźródłową (podrozdz.~\ref{sec:warianty:supersource}).
Algorytm uruchamiany jest nie z~pojedynczego wierzchołka, lecz z~superźródła
$s^\star_k$ połączonego krawędziami zerowej wagi ze~wszystkimi usługami
kategorii, więc do~izochrony trafia każdy wierzchołek znajdujący się
w~odległości $\le t$ od~\emph{którejkolwiek} usługi z~$S_k$. Dla realistycznych
miast, w~których liczność $|S_k|$ dochodzi do~kilkunastu tysięcy
(tab.~\ref{tab:uslugi_kategorie}), zbiór $C_k(t)$ może pokrywać dominującą
część grafu, a~praktyczny zysk z~samego ograniczenia progiem zależy w~znacznym
stopniu od~liczności i~rozkładu usług w~analizowanym obszarze. W~rozdziale~\ref{chap:wyniki}
zestawiono wyniki eksperymentów, które pozwalają oszacować ten zysk w~liczbach.

\subsection{Relacja do literatury}
\label{subsec:warianty:modified:relacja}

Odpowiada to algorytmowi \texttt{custom\_dijkstra} z~implementacji
referencyjnej Lama~\cite{Lam2024}. Różnice względem kodu użytego tutaj mają
charakter inżynieryjny i~dotyczą m.in.\
przygotowanie reprezentacji grafu (jednorazowe scalenie krawędzi równoległych
i~przeindeksowanie wierzchołków na liczby całkowite), wykorzystanie tablicowej
listy sąsiedztwa zamiast wielowarstwowych słowników biblioteki NetworkX
oraz dołączenie superźródła jako kolejnego wierzchołka w~tablicowej liście sąsiedztwa
(podrozdz.~\ref{sec:warianty:supersource}).

\section{Wariant z~kolejką kubełkową}
\label{sec:warianty:bucket}

W~podrozdziale~\ref{sec:intro:warianty-dijkstry} wskazano, że przy
ograniczonym zakresie odległości i~wąskim rozkładzie wag krawędzi możliwe
jest zastąpienie kopca binarnego \emph{kolejką kubełkową} (strukturą
zaproponowaną pierwotnie przez Diala~\cite{dial1969} i~uogólnioną przez
Cherkassky'ego, Goldberga i~Radzika~\cite{cherkassky1996}). Realizuje on tę samą logikę co~$t$-ograniczony algorytm Dijkstry, przy czym zamiast
kopca binarnego stosuje się jednopoziomową tablicę kubełków FIFO o~stałej
szerokości~$\Delta$. Wariant zaimplementowany w~niniejszej pracy
(podrozdz.~\ref{subsec:warianty:bucket:schemat}) odbiega od~klasycznego
opisu z~pracy~\cite{cherkassky1996} w~dwóch szczegółach: po~pierwsze,
pracujemy z~wagami rzeczywistymi (czas przejścia w~minutach), podczas
gdy analiza CGR zakłada wagi całkowite; po~drugie, zamiast jawnej
operacji \emph{decrease-key} polegającej na~przeniesieniu wierzchołka
między kubełkami przy zmianie etykiety stosujemy leniwe usuwanie -- ten
sam wierzchołek może występować w~kilku kopiach w~tym samym kubełku,
a~nieaktualne pary $(d_t,u)$ są odrzucane dopiero przy ekstrakcji.

\subsection{Struktura danych}
\label{subsec:warianty:bucket:struktura}

Kolejka kubełkowa składa się z~tablicy kubełków
$B[0], B[1],\dots, B[m-1]$, gdzie każdy kubełek jest kolejką FIFO
przechowującą pary $(d,v)$ (tymczasowa odległość $d$, wierzchołek $v$).
Wstawienie wierzchołka $v$ o~tymczasowej odległości $d$ polega na dopisaniu
pary $(d,v)$ na koniec kubełka o~indeksie
\[
\mathrm{idx}(d)=\lfloor d/\Delta\rfloor.
\]
Operacja \emph{extract-min} polega na liniowym skanowaniu tablicy
od~bieżącego wskaźnika \texttt{min\_idx} w~stronę większych indeksów do
chwili napotkania pierwszego niepustego kubełka i~zwróceniu pierwszej pary
$(d,v)$ z~tego kubełka.
Wskaźnik \texttt{min\_idx} nigdy się nie zmniejsza, dlatego łączna
praca wykonywana w~ramach skanu w~trakcie całego przebiegu algorytmu
ograniczona jest przez liczbę pól tablicy kubełków
\[
m=\lfloor t/\Delta\rfloor+1.
\]
W~implementacji przed każdym przeszukiwaniem długość tablicy kubełków jest
ustawiana tak, by wystarczyło co najmniej \(\lfloor t/\Delta\rfloor+1\) pól
(na indeksy \(0,\dots,\lfloor t/\Delta\rfloor\); w~szczególności dla
$t=15$~min i~$\Delta=1$~min jest to $16$~pól, bo odległość równa progowi
wpada do kubełka o~indeksie~$\lfloor t/\Delta\rfloor$).

\subsection{Schemat algorytmu}
\label{subsec:warianty:bucket:schemat}

\begin{codebox}
\Procname{$\proc{Bucket-Dijkstra}(G^\star_k, s^\star_k, t, \Delta)$}
\li \For each $v\in V\cup\{s^\star_k\}$
\li \Do $\id{dist}[v]\gets\infty$ \End
\li $\id{dist}[s^\star_k]\gets 0$
\li $\proc{Bucket-Insert}(B, s^\star_k, 0)$
\li \While $B$ zawiera niepusty kubełek
\li \Do $(d_t,u)\gets\proc{Bucket-Extract-Min}(B)$
\li     \If $d_t \neq \id{dist}[u]$
\li     \Then \kw{continue}                            \quad // nieaktualna para $(d_t,u)$
        \End
\li     \If $\id{dist}[u]\ge t$
\li     \Then \kw{continue}
        \End
\li     \For each $(u,v)\in E(G^\star_k)$
\li     \Do $\id{nd}\gets \id{dist}[u] + w(u,v)$
\li         \If $\id{nd}<\id{dist}[v]$ \kw{and} $\id{nd}\le t$
\li         \Then $\id{dist}[v]\gets\id{nd}$
\li             $\proc{Bucket-Insert}(B,v,\id{nd})$
            \End
        \End
    \End
\li \Return $\{v:\,\id{dist}[v]\le t\}$
\end{codebox}

Od wariantu ograniczonego progiem $t$ różni się więc głównie typem kolejki
oraz sposobem odrzucania nieaktualnych par: przy kopcu binarnym wystarcza
warunek $d>\id{dist}[u]$, a~przy kubełkach konieczne jest $d_t\neq\id{dist}[u]$,
bo przy $\Delta>w_{\min}$ ten sam wierzchołek może w~jednym kubełku czekać
w~kilku kopiach o~różnych tymczasowych odległościach (zob.\ podrozdz.\
\ref{subsec:warianty:bucket:tryby}).

\subsection{Tryby \emph{label-setting} i~\emph{label-correcting}}
\label{subsec:warianty:bucket:tryby}

W~kolejce kubełkowej liczy się stosunek szerokości
kubełka~$\Delta$ do najmniejszej dodatniej wagi krawędzi
$w_{\min}=\min_{e\in E}w(e)$. Od tego zależy, czy mamy do czynienia z~trybem
\emph{label-setting}, jak w~klasycznym algorytmie
Dijkstry, czy z~trybem \emph{label-correcting},
gdzie ten sam wierzchołek może być wyciągany z~kolejki wielokrotnie.
Podział na te dwa tryby i~odpowiadająca mu terminologia są standardowe
w~literaturze o~optymalizacji sieciowej; przyjmuję je za
Ahują, Magnantim i~Orlinem~\cite{ahuja1993network}.
Zachowanie wariantu kubełkowego o~szerokości~$\Delta$ jako algorytmu,
w~którym wierzchołki mogą być skanowane więcej niż raz, odnotowali
eksperymentalnie Cherkassky, Goldberg i~Radzik przy okazji analizy
\emph{approximate bucket implementation} algorytmu
Dijkstry~\cite{cherkassky1996}.

\paragraph{\emph{label-setting} ($\Delta\le w_{\min}$).}
Przy $\Delta\le w_{\min}$ wariant kubełkowy zachowuje kluczową własność
klasycznego algorytmu Dijkstry: każdy wierzchołek jest wyciągany z~kolejki
co~najwyżej raz, a~wartość $\id{dist}[u]$ w~momencie wyciągnięcia jest już
docelowa. Argument przebiega standardowo: gdy procedura wyciąga wierzchołek
$u$ z~kubełka o~indeksie $k$, wszystkie jeszcze nieobsłużone wierzchołki $v$
spełniają $\id{dist}[v]\ge k\Delta$ (wskaźnik \texttt{min\_idx} nigdy się nie
cofa). Każda ścieżka, która mogłaby jeszcze obniżyć $\id{dist}[u]$, musiałaby kończyć
się krawędzią wychodzącą z~pewnego takiego $v$ i~prowadzącą do~$u$, dając
wartość
\[
\id{dist}[v]+w(v,u)\;\ge\;k\Delta+w_{\min}\;\ge\;k\Delta+\Delta\;=\;(k+1)\Delta,
\]
czyli wyprowadzającą $\id{dist}[u]$ poza bieżący kubełek. W~chwili wyciągnięcia
$\id{dist}[u]$ jest więc już ostateczne. Ten sam argument leży u~podstaw analizy
Cherkassky'ego, Goldberga i~Radzika~\cite{cherkassky1996}; algorytm Diala dla
wag całkowitych jest jego szczególnym przypadkiem, w~którym $\Delta=1$ oraz
$w_{\min}=1$~\cite{dial1969}.

\paragraph{\emph{label-correcting} ($\Delta>w_{\min}$).}
Gdy $\Delta>w_{\min}$, w~obrębie jednego kubełka mogą znaleźć się wierzchołki,
których tymczasowe odległości różnią się więcej niż o~$w_{\min}$. Wówczas
relaksacja krawędzi $(u,v)$, w~której zarówno $u$, jak i~$v$ leżą w~tym samym
kubełku, może obniżyć $\id{dist}[v]$ i~wstawić $v$ do kolejki kolejny raz,
być może ponownie do tego samego kubełka. W~konsekwencji ten sam wierzchołek
trafia do kolejki wielokrotnie, a~powyższy argument przestaje obowiązywać:
$\id{dist}[u]$ w~chwili pierwszego wyciągnięcia nie musi być jeszcze
ostateczne.

\clearpage

Poprawność zachowuje mechanizm \emph{lazy deletion}: każda para $(d_t,u)$
przechowuje wartość $\id{dist}[u]$ z~chwili wstawienia, a~przy wyciąganiu
sprawdzamy warunek $d_t=\id{dist}[u]$. Pary nieaktualne (dla których
$d_t>\id{dist}[u]$, bo wartość w~tablicy została w~międzyczasie poprawiona)
są pomijane bez relaksacji. Każda \emph{skuteczna} ekstrakcja, czyli taka,
która faktycznie wywołuje relaksację sąsiadów, odpowiada nowej, mniejszej
wartości $\id{dist}[u]$ niż jakakolwiek wcześniejsza skuteczna ekstrakcja
tego wierzchołka. Ciąg kolejnych wartości $\id{dist}[u]$ jest więc
malejący; ponieważ każda z~tych wartości jest długością pewnej ścieżki
z~$s^\star_k$ do~$u$ w~skończonym grafie o~nieujemnych wagach, zbiór możliwych
wartości jest skończony, więc ciąg nie może ciągnąć się w~nieskończoność.
Algorytm wykonuje zatem skończoną liczbę kroków. Wielkość tego nadmiaru
względem trybu \emph{label-setting} zależy od rozkładu wag krawędzi i~rozrzutu
odległości w~obrębie pojedynczych kubełków; omówieniu tego efektu poświęcony
jest podrozdz.~\ref{subsec:warianty:bucket:delta}.

\subsection{Dobór szerokości kubełka $\Delta$}
\label{subsec:warianty:bucket:delta}

Wybór $\Delta$ jest kluczowym parametrem wariantu kubełkowego.
W~eksperymentach przyjęto domyślnie $\Delta=1{,}0$~min.
Uzasadnienie wynika z~właściwości użytych grafów pieszych, a~nie z~chęci
wymuszenia $\Delta\le w_{\min}$ i~trybu \emph{label-setting}.

Empirycznie, na grafach pieszych Gdańska, Gdańska Południe i~Londynu
zachodzą następujące zależności:
\begin{itemize}
  \item mediana wagi krawędzi pieszej wynosi ok.\ $0{,}7$~min,
        co~odpowiada odcinkowi długości ok.\ $58$~m przy prędkości
        $5$~km/h (podrozdz.~\ref{sec:dane:wagi}),
  \item minimalna dodatnia waga krawędzi wynosi ok.\ $7\cdot 10^{-4}$~min;
        wartość ta jest zdominowana przez krótkie segmenty powstałe
        w~procedurze wstawiania usług do grafu (rozdz.~\ref{sec:dane:wstawianie}),
        gdy wstawiany wierzchołek usługi pada blisko końca krawędzi.
\end{itemize}

Wybór $\Delta=w_{\min}$ formalnie zapewniłby tryb \emph{label-setting}, jednak
przy progu $t=15$~min skutkowałby liczbą kubełków
$\lceil t/w_{\min}\rceil\approx 21\,500$. W~każdym z~nich znajdowałaby się
średnio mniej niż jedna para $(d,v)$, a~koszt operacji \emph{extract-min} byłby
zdominowany przez liniowy skan tablicy pustych kubełków. Asymptotyczna
przewaga $O(1)$ na operację (właściwa dla wariantu kubełkowego) traciłaby
wówczas znaczenie praktyczne, a~mierzony czas wykonania okazałby się
gorszy niż dla wariantu z~kopcem binarnym.

\clearpage

Wybór $\Delta=1{,}0$ daje natomiast:
\begin{itemize}
  \item liczbę kubełków $\lceil t/\Delta\rceil=15$, czyli stałą znacznie
        mniejszą niż liczba wierzchołków w~izochronie,
  \item gęste obsadzenie pojedynczego kubełka: przy $|V'|\gg\lceil t/\Delta\rceil=15$
        (tab.~\ref{tab:wyniki:pokrycie}) na jeden kubełek przypada
        od kilkuset do dziesiątek tysięcy par $(d,v)$, więc koszt skanu
        przy \emph{extract-min} dotyczy krótkiej, stale zajętej tablicy,
        a~nie długiej sekwencji pustych kubełków,
  \item dopuszczalność trybu \emph{label-correcting}; ze~względu na wąski rozkład
        wag krawędzi liczba powtórnych ekstrakcji pojedynczego wierzchołka
        pozostaje w~praktyce niewielka, a~jej koszt jest amortyzowany przez
        bardzo niski koszt operacji na kolejce.
\end{itemize}

W~niniejszej pracy nie zbadano, czy przyjęta wartość $\Delta=1{,}0$~min
jest optymalnym wyborem tego parametru.

\subsection{Złożoność}
\label{subsec:warianty:bucket:zlozonosc}

W~trybie \emph{label-setting} ($\Delta\le w_{\min}$) wariant kubełkowy daje
amortyzowany koszt $O(1)$ na operację \emph{insert} oraz \emph{extract-min},
plus jednorazowy skan tablicy o~koszcie $O(\lceil t/\Delta\rceil)$.
Łączna złożoność wynosi wówczas
\[
T_{\mathrm{bucket}}=\mathcal{O}\!\bigl(|V'|+|E'|+\lceil t/\Delta\rceil\bigr),
\]
gdzie, jak poprzednio, $V'$ i~$E'$ oznaczają część grafu faktycznie odwiedzoną
przy ograniczeniu progiem~$t$. W~porównaniu z~kopcem binarnym
($\mathcal{O}((|V'|+|E'|)\log|V'|)$) usunięty zostaje czynnik logarytmiczny.

W~trybie \emph{label-correcting} ($\Delta>w_{\min}$) pojedynczy wierzchołek
bywa wyciągany z~kolejki więcej niż raz; w~praktyce, przy wąskim rozkładzie
wag na sieci pieszej, liczba takich powtórzeń i~tak pozostaje mała.
Asymptotyczny rząd kosztu nie zmienia się względem \emph{label-setting},
rosną jedynie stałe ukryte w~notacji $\mathcal{O}(\cdot)$.

\clearpage

\section{Podsumowanie i porównanie wariantów}
\label{sec:warianty:podsumowanie}

Tabela~\ref{tab:warianty_porownanie} zbiera w~skrócie trzy warianty: rodzaj
kolejki, stosowanie progu~$t$ oraz porządek złożoności czasowej w~ujęciu
istotnym dla zadania miasta 15-minutowego.

\begin{table}[ht]
\centering
\caption{Porównanie trzech rozważanych wariantów algorytmu Dijkstry.
$V'$ i~$E'$ oznaczają część grafu faktycznie odwiedzaną przy ograniczeniu
progiem~$t$.}
\label{tab:warianty_porownanie}
\begin{tabular}{l l l l}
\hline
\textbf{Wariant} & \textbf{Rodzaj kolejki} & \textbf{Ograniczenie $t$} & \textbf{Złożoność czasowa} \\
\hline
bazowy           & kopiec binarny    & Nie & $\mathcal{O}((|V|+|E|)\log|V|)$            \\
ograniczony~$t$  & kopiec binarny    & Tak          & $\mathcal{O}((|V'|+|E'|)\log|V'|)$         \\
kubełkowy        & tablica kubełków FIFO    & Tak          & $\mathcal{O}(|V'|+|E'|+\lceil t/\Delta\rceil)$ \\
\hline
\end{tabular}
\end{table}

Trzy warianty stanowią naturalną sekwencję rozszerzeń: wariant bazowy nie
wykorzystuje ani ograniczonego zasięgu izochrony, ani struktury rozkładu
wag krawędzi; wariant ograniczony progiem~$t$ uwzględnia tę pierwszą cechę,
nie zmieniając kolejki priorytetowej; wariant kubełkowy łączy obie obserwacje
i~dodatkowo redukuje koszt operacji na kolejce.

W~rozdziale~\ref{chap:wyniki} zestawimy te szacunki z~pomiarami czasu na
grafach pieszych Gdańska, Gdańska Południe i~Londynu. Wcześniej,
w~rozdziale~\ref{chap:implementacja}, opisany jest wspólny kod trzech wariantów
oraz tryb wieloprocesowy z~równoległym liczeniem kategorii usług na osobnych
wątkach.
